\section{Introduction}
\label{sec:intro}

Congressional redistricting in the United States requires satisfying multiple
simultaneous criteria: equal population (the one-person-one-vote standard),
geographic compactness, contiguity, and preservation of political subdivisions
such as counties, cities, and townships.
While the first three criteria have been extensively studied in the algorithmic
redistricting literature~\citep{Duchin2019, Mehrotra1998}, the fourth---preserving
political subdivisions---has received comparatively little formal treatment.

This omission is legally significant.
Most state redistricting constitutions and statutes list subdivision preservation
as an explicit criterion.
California's constitution requires that ``the geographic integrity of any city,
county, or city and county be respected.''
Florida requires that districts not ``be drawn to favor or disfavor an incumbent
or political party.''
The North Carolina Supreme Court in \emph{Harper v.\ Hall} (2022) invalidated
a map in part for excessive county splits.
Yet when redistricting commissions adopt algorithmic methods, no existing system
provides a formal mechanism for trading off subdivision preservation against
compactness and population balance.

We fill this gap.
The key insight is that census tract GEOIDs encode county membership in their
first five digits (two-digit state FIPS plus three-digit county FIPS within state),
so county affiliation is freely available from TIGER shapefiles with no additional
data acquisition.
We introduce a \emph{stickiness parameter} $\alpha \geq 0$ that scales the weight
of intra-county edges in the METIS graph:
\begin{equation}
  w'(u,v) = w(u,v) \cdot \bigl(1 + \alpha \cdot \mathbf{1}[\text{county}(u) = \text{county}(v)]\bigr),
  \label{eq:weight}
\end{equation}
where $w(u,v)$ is the TIGER boundary length between tracts $u$ and $v$.
Setting $\alpha = 0$ recovers standard minimum-edge-cut (MEC) bisection.
Setting $\alpha > 0$ makes the algorithm prefer cuts that run along existing
county lines, since those edges have lower (unmodified) weight.

Our 50-state empirical evaluation (Congressional districts, 2020 Census, $\alpha = 5$)
shows county splits fall by 45\% on average relative to MEC, with no systematic
partisan bias.
The tradeoff is a 2.5$\times$ increase in total boundary length---the algorithm
is drawing longer but more politically meaningful cuts.

The stickiness parameter is a single published number.
Any independent analyst with access to the Census TIGER files can verify that
a given plan was produced by running METIS with $\alpha = 5$, confirming that
county splits are not a product of partisan choice but of the mechanical
application of a statutory requirement.

\paragraph{Contributions.}
\begin{enumerate}
  \item We formalize subdivision preservation as a tunable edge-weight signal
    (Section~\ref{sec:method}), show it is derivable from TIGER data alone
    (no partisan inputs), and prove it reduces to MEC when $\alpha = 0$.
  \item We provide a 50-state empirical evaluation (Section~\ref{sec:eval})
    measuring county splits, edge cut, and partisan outcomes at $\alpha \in \{0, 5\}$.
  \item We demonstrate that subdivision-respecting redistricting is compositional:
    $\alpha$ can be combined with any split strategy (bisection, GeoSection,
    CompactBisect) and any other edge-weight signal (partisan, minority) through
    the $\mathtt{EdgeWeighter}$ pipeline introduced in this series.
\end{enumerate}
